% Generated draft: requires mathematical and source review.
\documentclass[11pt]{article}
\usepackage[margin=25mm]{geometry}
\usepackage{fontspec}
\setmainfont{Latin Modern Roman}
\usepackage{amsmath,amssymb,mathtools,unicode-math}
\setmathfont{Latin Modern Math}
\usepackage{enumitem,needspace,xurl,hyperref}
\setlength{\parindent}{0pt}
\setlength{\parskip}{5pt}
\setlength{\emergencystretch}{3em}
\newcommand{\N}{\mathbb{N}}
\newcommand{\Z}{\mathbb{Z}}
\newcommand{\Q}{\mathbb{Q}}
\newcommand{\R}{\mathbb{R}}
\newcommand{\C}{\mathbb{C}}
\newcommand{\F}{\mathbb{F}}
\newcommand{\A}{\mathbb{A}}
\newcommand{\PP}{\mathbb{P}}
\newcommand{\E}{\mathbb{E}}
\newcommand{\eps}{\varepsilon}
\newcommand{\dd}{\,\mathrm{d}}
\begin{document}
\section*{UnsolvedMath ID: 30003233}
\textbf{Source problem code:} catalogue entry
\subsection{Definitions and mathematical statement}
For a real number \(x\), let \(\|x\|:=\min_{n\in\Z}|x-n|\) denote its distance to the nearest integer. Interpret \(q\to\infty\) as tending to infinity through positive integers, and take \(\log\) to be the natural logarithm. For each real parameter \(\lambda\), define
\[
M_\lambda:=\left\{(\alpha,\beta)\in\R^2:\liminf_{\substack{q\to\infty\\q\in\N,\ q\ge2}}q(\log q)^\lambda\|q\alpha\|\,\|q\beta\|>0\right\}.
\]
Here \(\dim_{\mathrm H}\) denotes Hausdorff dimension in the Euclidean plane, so full Hausdorff dimension means dimension \(2\).

The supplied problem asks for a correct multiplicative analogue of badly approximable points and, specifically, asks whether both assertions hold:
\[
\begin{aligned}
&\forall\lambda\in\R,\quad \lambda<1\ \Longrightarrow\ M_\lambda=\varnothing,\\
&\forall\lambda\in\R,\quad \lambda\ge1\ \Longrightarrow\ \dim_{\mathrm H}(M_\lambda)=2.
\end{aligned}
\]
The boundary value \(\lambda=1\) is included in the second assertion. These are questions, not asserted conclusions. The integer domain for \(q\), the real domain for \(\lambda\), and the meaning of \(\|\cdot\|\) are conventional interpretations of the supplied notation; they are not explicitly defined in the source record. The broader request for a ``correct multiplicative analogue'' supplies no separate criterion beyond these concrete questions.
\subsection{Short English statement}
For each real \(\lambda\), consider pairs of real numbers for which the lower limit of \(q(\log q)^\lambda\|q\alpha\|\|q\beta\|\), over positive integers \(q\), is strictly positive. Must there be no such pairs when \(\lambda<1\), and must their set have Hausdorff dimension \(2\) when \(\lambda\ge1\)? The supplied record labels this problem open; that status has not been independently verified.
\subsection{Sources}
\begin{itemize}
\item Local UnsolvedMath registry, permanent ID \(30003233\), title ``Multiplicatively Badly Approximable Pairs,'' source problem code \texttt{catalogue entry}, in the collection ``Oberwolfach Reports Open Problems.'' No individual proposer is supplied; the proposed year is recorded as 2016.
\item Supplied report citation: \emph{Arbeitsgemeinschaft: Diophantine Approximation, Fractal Geometry and Dynamics} (2017), Oberwolfach Reports, pp.\ 2749--2792, DOI \texttt{10.4171/owr/2016/48}. \url{https://doi.org/10.4171/owr/2016/48}. The supplied background describes these reports as documenting mathematical workshops and their open-problem discussions. The relationship between the citation's 2017 date and the record's 2016 proposed year and DOI year is not explained.
\item Supplied literature source: Badziahin, \emph{On multiplicatively badly approximable numbers}. \url{https://arxiv.org/abs/1101.1855}.
\item Supplied literature source: ``Oberwolfach Report 48.'' \url{https://doi.org/10.4171/owr/2016/48}.
\item Literature assessment dated exactly \texttt{2026-08-21}: the supplied curation labels the problem open, relates multiplicatively badly approximable pairs to Littlewood's conjecture, described there as still open, and states that logarithmic-factor constructions yield large exceptional sets without answering the zero-factor problem. This is supplied dated curation derived from the OWR material, not an independently verified assessment. Its time-sensitive status and cited claims require independent checking.
\end{itemize}
\end{document}
